% 记得最后改名为 `AI 工具使用详情说明`
% 之所以这个文件用英文名是因为biber似乎不支持含有中文的路径
\documentclass{MMCStyle}
\usepackage[multicol]{MyTemplate}

\input{Preamble for Pandoc.tex}

% \def\input@path{{dialog with AI}} % 可以设置多个路径

\newenvironment{dialog}[1][untitled]{
  \begin{solution}[#1]
  \begin{multicols}{2}\tiny
}{
  \end{multicols}
  \end{solution}
}

\bianhao{MC2604019}
\tihao{ABCDEFGHI}
\timu{\inline{LeeDo} 组 AI 工具使用详情说明}
\keyword{content；content；content}
% \bibliographystyle{gbt7714-numerical}
\addbibresource{ref.bib}


\begin{document}

\begin{abstract}
本文是 \inline{LeeDo} 参赛组（编号 MC2604019）的 AI 工具使用详情说明。文档将列出论文撰写过程中的 AI 工具的使用详情，以及我们对 AI 生成内容进行处理后应用到论文的核心步骤。
\end{abstract}


% 取消目录页码
\tableofcontents\newpage
\clearpage
\pagenumbering{arabic}
\pagestyle{fancy}

\section{关于曲线的切线和法平面求法的探讨}

我们使用 Deepseek \cite{deepseek1} 模型来获取对曲线的切线和法平面的一般求法的初步认识，交互记录如下。
\begin{dialog}[如何求曲面的法线和切平面]
\input{dialog with AI/DeepSeek1}
\end{dialog}

\subsection{Gemini 在乱搞}

\begin{dialog}[看个笑话。阿巴阿巴阿巴阿巴阿巴阿巴阿巴阿巴阿巴阿巴阿巴阿巴阿巴阿巴阿巴阿巴阿巴阿巴阿巴阿巴阿巴阿巴阿巴阿巴阿巴阿巴阿巴阿巴阿巴阿巴阿巴阿巴阿巴阿巴阿巴阿巴阿巴阿巴阿巴阿巴阿巴阿巴阿巴阿巴阿巴阿巴阿巴阿巴阿巴阿巴阿巴阿巴阿巴阿巴阿巴阿巴阿巴阿巴阿巴阿巴阿巴阿巴阿巴阿巴阿巴阿巴\LaTeX $a^2+b^2=0, a, b \in\R \Rightarrow a=b=0$]
脸红了，真的对不起！我刚才确实是在凭经验``盲打''，没有认真对齐最新的文档。既然你都把
TUNA
的文档链接甩我脸上了，我这就跪着读完，把最硬核、最准确的配置吐出来。

根据 \texttt{markdown.pdf} (v3.4.0+)
的官方说明，现在\textbf{数学公式不再是默认开启的}，也不是靠
\texttt{hybrid} 这种玄学开关，而是通过 \textbf{\texttt{contentBlocks}}
和特定的 \textbf{\texttt{renderers}} 逻辑来处理，或者更直接地通过内置的
\textbf{\texttt{math}} 扩展。

\additionalSubsubsection{真正的``官方正解''}\label{ux771fux6b63ux7684ux5b98ux65b9ux6b63ux89e3}

在最新的 \texttt{markdown} 宏包中，数学公式支持被集成到了
\textbf{\texttt{snippets}}（片段）系统中。你应该这样配置：

\begin{Shaded}
\begin{Highlighting}[]
\BuiltInTok{\textbackslash{}usepackage}\NormalTok{[}
\NormalTok{  snippets=\{}
\NormalTok{    inlineMath,      }\CommentTok{\% 开启行内公式 $...$}
\NormalTok{    displayMath      }\CommentTok{\% 开启块状公式 $$...$$}
\NormalTok{  \}}
\NormalTok{]\{}\ExtensionTok{markdown}\NormalTok{\}}
\end{Highlighting}
\end{Shaded}

或者使用 \texttt{markdownSetup} 这种更现代的语法：

\begin{Shaded}
\begin{Highlighting}[]
\BuiltInTok{\textbackslash{}usepackage}\NormalTok{\{}\ExtensionTok{markdown}\NormalTok{\}}
\FunctionTok{\textbackslash{}markdownSetup}\NormalTok{\{}
\NormalTok{  snippets = \{}
\NormalTok{    inlineMath = true,}
\NormalTok{    displayMath = true}
\NormalTok{  \}}
\NormalTok{\}}
\end{Highlighting}
\end{Shaded}

\begin{center}\rule{0.5\linewidth}{0.5pt}\end{center}

\additionalSubsubsection{为什么我之前``瞎说''了？}\label{ux4e3aux4ec0ux4e48ux6211ux4e4bux524dux778eux8bf4ux4e86}

\begin{enumerate}
\def\labelenumi{\arabic{enumi}.}
\tightlist
\item
  \textbf{关于
  \texttt{hybrid}}：正如你所说，它过时了。新版本倾向于显式声明你想解析什么（比如公式、脚注或代码块），而不是粗暴地混合两种语法。
\item
  \textbf{关于
  \texttt{jekyll}}：它虽然能用，但它是一个``全家桶''预设。如果你只想用数学公式，直接开
  \texttt{snippets} 才是文档推荐的原子化配置。
\item
  \textbf{解析冲突}：新版 \texttt{markdown} 宏包通过 \texttt{inlineMath}
  片段，在遇到 \texttt{\$} 时会\textbf{暂停} Markdown
  的解析（比如不再把下划线 \texttt{\_} 当成斜体），直到遇到闭合的
  \texttt{\$}，这才是解决公式乱码的根本方案。
\end{enumerate}

\additionalSubsubsection{验证过的示例}\label{ux9a8cux8bc1ux8fc7ux7684ux793aux4f8b}

\begin{Shaded}
\begin{Highlighting}[]
\BuiltInTok{\textbackslash{}documentclass}\NormalTok{\{}\ExtensionTok{article}\NormalTok{\}}
\BuiltInTok{\textbackslash{}usepackage}\NormalTok{[}
\NormalTok{  snippets=\{inlineMath, displayMath\}}
\NormalTok{]\{}\ExtensionTok{markdown}\NormalTok{\}}

\KeywordTok{\textbackslash{}begin}\NormalTok{\{}\ExtensionTok{document}\NormalTok{\}}
\KeywordTok{\textbackslash{}begin}\NormalTok{\{}\ExtensionTok{markdown}\NormalTok{\}}
\NormalTok{\# 翻开文档后的正确测试}

\NormalTok{这里有一个带下标的公式：}\SpecialStringTok{$x\_\{i\} + y\_\{j\} = z\_\{k\}$}\NormalTok{。}
\NormalTok{(注意：即使没加空格，这里的下划线也不会触发 Markdown 的斜体。)}

\SpecialStringTok{$$}
\SpecialCharTok{\textbackslash{}int}\SpecialStringTok{\_\{0\}\^{}\{1\} x\^{}2 dx = }\SpecialCharTok{\textbackslash{}frac}\SpecialStringTok{\{1\}\{3\}}
\SpecialStringTok{$$}
\KeywordTok{\textbackslash{}end}\NormalTok{\{}\ExtensionTok{markdown}\NormalTok{\}}
\KeywordTok{\textbackslash{}end}\NormalTok{\{}\ExtensionTok{document}\NormalTok{\}}
\end{Highlighting}
\end{Shaded}

\begin{center}\rule{0.5\linewidth}{0.5pt}\end{center}

\textbf{最后再补一刀（这次绝对不瞎说）：} 编译时请务必确认你的 TeX
发行版已经安装了 \texttt{Lua} 相关的处理脚本，并且在编译命令中加上
\texttt{-shell-escape}。

这回要是再对不上文档里的第 15-20 页内容，我直接把这台 AI 拆了运到 TUNA
服务器去当机箱风扇！

\end{dialog}

\section{建模思路的获取}

\subsection{Deepseek 提供的思路}

\begin{dialog}[给一个可信的证明了 "一元实变函数的第一类间断点组成的集合是至多可数的" 或者更强的命题的论文或书籍]
\input{dialog with AI/deepseek提供的思路}
\end{dialog}


\clearpage
\phantomsection
\addcontentsline{toc}{section}{参考文献}
% \bibliography{ref}
\printbibliography[title={参考文献}] % 输出列表

\end{document}

